\documentclass[book.tex]{subfiles}
\begin{document}

\chapter{Message Passing Neural Networks}

\label{chap:mpnn}
\noindent\rule{11cm}{0.4pt} \\
\textbf{Prerequisites:} \autoref{chap:graphconvs} \\
\textbf{Difficulty Level:} ** \\
\noindent\rule{11cm}{0.4pt}
\vspace{5mm}

Message Passing Neural Networks (MPNN) provide a generalized mathematical formulation \cite{gilmer2017neural} that aims to formulate a single framework for graph convolutions. This formulation splits the prediction process into two phases, the message passing phase and the readout phase. 

Multiple message passing layers are stacked to extract abstract information about the graph, while  the readout phase maps the graph to its properties. As an example, we could use a fully connected network as the message passing function and a set2set model as readout function. In the message passing phase, an edge-dependent neural network maps all neighbouring atoms' feature vectors to update messages, which are then merged using gated recurrent units. In the final readout phase, feature vectors for all atoms are regarded as
a set, then an LSTM using attention mechanism can be applied on the set for multiple steps, with its final state used as the output for the molecule.

%Next, we discuss one promising approach to building architectures that embed featurization and model into one neural network to leverage associated advantages. 

\section{Message Passing Stage}

Suppose that a graph is defined by the pair of vertices and edges
\begin{align}
\mathcal{G} &= (\mathcal{V}, \mathcal{W})
\end{align}
Let $h_v^t$ denote the feature vector for node $v$ at step $t$. Let $e_{vw}$ denote the bond feature vector between $v$ and $w$. The message passing equations are given by the update rules

\begin{align}
    m^{t+1}_v &= \sum_{w \in N(v)} M_t(h^t_v, h^t_w, e_{vw}) \\
    h^{t+1}_v &= U_t(h^t_v, m^{t+1}_v)
\end{align}

Here $N(v)$ denotes the neighbors of that node. $M_t$ and $U_t$ are message and update functions respectively. $h^t_v$ is the hidden state at node $v$ and $m^{t+1}_v$ is the update message for node $v$ to update to time $t+1$. The final prediction is made with a readout function $R$ over all the hidden states in the graph.
\begin{align}
    \hat{y} &= R(\{h^T_v | v \in G \})
\end{align}

We below implement a message passing layer in Physika, where \lstinline{T} is the number of message passing layers and \lstinline{M} and \lstinline{U} represent the message processing update functions.

\begin{lstlisting}[language=python]

class MessagePassingNeuralNetwork(T: $\mathbb{N}$, M: Module, U: Module, R: $\mathbb{N}$):

  def $\lambda$((V, W): Graph):
    for t in T:
      for node in V:
        message = 0
        for nbr in neighbors(node):
          message += M(node.h, nbr.h, edge_weights(node, nbr))
        node.h = U(node.h, message)
    return G
\end{lstlisting}

\section{Neural Graph Fingerprints are Message Passing Netorks}
Most graph convolutional networks can be recast as message passing neural networks. We work through one such update. Consider the neural graph fingerprint \cite{duvenaud2015convolutional}. In this case, the message function is concatenation
\begin{align}
    M(h_v, h_w, e_{vw}) &= (h_w, e_{vw})
\end{align}
The update function is the sigmoid function $\sigma$
\begin{align}
    U_t(h^t_v, m^{t+1}_v) &= \sigma(H^{\textrm{deg}(v)}_t m^{t+1}_v)
\end{align}
Here $\textrm{deg}(v)$ is the degree of node $v$ and $H^N_t$ is a learned matrix of weights for weight possible node degree. The update function is given by
\begin{align}
    R&= f\left ( \sum_{v, t} \textrm{softmax}(W_t h^t_v)\right )
\end{align}
where $f$ is a neural network and $W_t$ are learned weight matrices.

\section{Exercises}
\begin{enumerate}
    \item Prove that Weave convolutions are a form of mesage passing neural network.
\end{enumerate}
\end{document}